% Chapter file: 023_Bell_De_ch.tex
% Source: 023_Bell_De.tex

% Original:	
\begin{tcolorbox}[colback=red!3!white,colframe=red!50!black,boxrule=0.8pt,title={Korrekturhinweise},fonttitle=\bfseries]
\textbf{Korrekturen aus der Rechenprüfung (1.~Okt.~2026).} Die folgenden Stellen sind im Text korrigiert bzw. vermerkt; jede trägt dort einen datierten Hinweis mit dem vorherigen Wert:
\begin{itemize}
\item CHSH$^{\mathrm{T0}}$: $2.827$ ($\Delta = 0.04$ \%) $\to 2.82805$ ($\Delta = 0.013$ \%); Aussagen \glqq Lokalität wiederhergestellt\grqq{}, \glqq ohne Bells Ungleichung zu verletzen\grqq{} und \glqq CHSH$^{\mathrm{ext}} < 2.5$\grqq{} $\to$ Bell-Verletzung bleibt erhalten ($\approx 2.828 > 2$); Faktor $K_{\mathrm{frak}}^{D_f}$ in Gl.~(\ref{eq:chsh_t0}) entfernt (ergäbe 2.717); Tabelle~\ref{tab:bell_ml}: Vorzeichen bei $\pi$ und $5\pi/4$ ($-1.000 \to +1.000$, $-0.707 \to +0.707$).
\end{itemize}
\end{tcolorbox}
	
	\section*{Abstract}
	Diese Erweiterung der T0-Serie wendet Erkenntnisse aus vorherigen \\ML-Tests (Wasserstoff-Niveaus) auf Bell-Tests an, um Quantenverschränkung im T0-Rahmen zu modellieren. Basierend auf der Zeit-Masse-Dualität und $\xi = 4/30000$ werden Korrelationen $E(a,b) = -\cos(a-b) \cdot (1 - \xi \cdot f(n,l,j))$ modifiziert, wobei $f(n,l,j)$ aus T0-Quantenzahlen stammt. Ein PyTorch-NN (1→32→16→1, 200 Epochen) simuliert CHSH-Verletzungen mit T0-Dämpfung, ergibt eine Reduktion von 2.82843 auf 2.82805 ($\Delta \approx 0.013$ \%); die Bell-Verletzung bleibt dabei praktisch vollständig erhalten (2.828 liegt 41 \% über der lokalen Schranke 2). \textit{(Korrigiert am 1.~Okt.~2026: vorher \glqq von 2.828 auf 2.827 (0.04 \% $\Delta$), was Lokalität bei $\xi$-Skala wiederherstellt\grqq{} -- $2\sqrt{2}\,(1-\xi) = 2.82843 \cdot (1 - 1.333 \times 10^{-4}) = 2.82805$, $\Delta = 0.013$ \%; 2.827 und 0.04 \% entstehen nur durch Abschneiden von $2\sqrt{2}$ auf 2.828. Jeder Wert $> 2$ verletzt CHSH; eine Dämpfung der Größe $10^{-4}$ stellt keine Lokalität her.)} Neue Erkenntnisse: ML zeigt subtile nicht-lokale Effekte als emergente Zeitfeld-Fluktuationen; Divergenz bei hohen Winkeln deutet auf fraktale Pfad-Interferenz hin. Bells Ungleichung bleibt dabei verletzt; die $\xi$-Dämpfung ist eine Korrektur der Größe $10^{-4}$ \textit{(Korrigiert am 1.~Okt.~2026: vorher \glqq Dies löst das EPR-Paradoxon harmonisch, ohne Bells Ungleichung zu verletzen\grqq{} -- mit CHSH$^{\mathrm{T0}} \approx 2.828 > 2$ ist die Ungleichung verletzt; lokale verborgene Variablen sind durch lückenlose Bell-Tests seit 2015 ausgeschlossen.)} – testbar via 2025-Loophole-free Experimente (z.\,B. 73-Qubit-Lie-Detector). Kaum Vorteile durch ML: Die harmonische T0-Berechnung ($\phi$-Skalierung) liefert bereits exakte Vorhersagen; ML kalibriert nur ($\sim$0.1 \% Genauigkeitsgewinn).
	
	\section{Einführung: Bell-Tests im T0-Kontext}
	\label{sec:intro_bell}
	
	Bell-Tests testen Quantenverschränkung vs. lokale Realität: Standard-QM verletzt Bells Ungleichung (CHSH >2), implizierend Nicht-Lokalität (EPR-Paradoxon). T0 löst dies durch $\xi$-modifizierte Korrelationen: Zeitfeld-Fluktuationen dämpfen die Korrelationen um einen Faktor der Größe $10^{-4}$; lokaler Realismus wird dadurch nicht hergestellt. \textit{(Korrigiert am 1.~Okt.~2026: vorher \glqq dämpfen Verschränkung lokal, bewahrend Realismus\grqq{} -- CHSH bleibt bei $\approx 2.828$, 41 \% über der lokalen Schranke 2.)} Basierend auf ML-Tests aus QM-Doc (Divergenz bei hohen $n$), simulieren wir hier CHSH mit T0-Korrekturen.
	
	\textbf{2025-Kontext:} Neueste Experimente (z.\,B. 73-Qubit-Lie-Detector, Oct 2025)\cite{sciencedaily2025} bestätigen QM-Verletzungen; T0 vorhersagt subtile Abweichungen ($\Delta \sim 10^{-4}$), testbar in Loophole-free Setups.
	
	Parameter: $\xi=4/30000$, $\phi \approx 1.618$; Quantenzahlen für Photonenpaare: $(n=1,l=0,j=1)$ (Photonen als Gen-1).
	
	\section{T0-Modifikation der Bell-Korrelationen}
	\label{sec:mod}
	
	Standard: $E(a,b) = -\cos(a-b)$ für Singulett-Zustand; CHSH = $E(a,b) - E(a,b') + E(a',b) + E(a',b') \approx 2\sqrt{2} \approx 2.828 >2$.
	
	T0: Zeitfeld dämpft: $E^{\mathrm{T0}}(a,b) = -\cos(a-b) \cdot (1 - \xi \cdot f(n,l,j))$, mit $f(n,l,j) = (n/\phi)^l \cdot [1 + \xi j / \pi] \approx 1$ (für Photonen). Dies reduziert CHSH auf $2\sqrt{2} \cdot (1 - \xi) \approx 2.82805$ ($\Delta \approx 0.013$ \%), also 41 \% über der lokalen Schranke 2 – die Bell-Verletzung bleibt erhalten. \textit{(Korrigiert am 1.~Okt.~2026: vorher \glqq $\approx 2.828 \cdot (1 - \xi) \approx 2.827$, knapp über 2 – Lokalität bei $\xi$-Präzision\grqq{} -- $2.82843 \cdot (1 - 1.333 \times 10^{-4}) = 2.82805$; $2.828/2 = 1.414$, eine Dämpfung um $10^{-4}$ kann den Wert nicht an 2 heranführen.)}
	
	\begin{equation}
		\mathrm{CHSH}^{\mathrm{T0}} = 2\sqrt{2} \cdot (1 - \xi \cdot \Delta \theta / \pi),
		\label{eq:chsh_t0}
	\end{equation}
	wobei $\Delta \theta = |a-b|$ (Winkelunterschied), $D_f=3-\xi$. \textit{(Korrigiert am 1.~Okt.~2026: vorher mit Faktor $K_{\mathrm{frak}}^{D_f}$ -- mit $K_{\mathrm{frak}} = 1 - 100\xi$ ist $K_{\mathrm{frak}}^{D_f} = 0.9605$ und damit CHSH$^{\mathrm{T0}} \approx 2.717$, was dem Zahlenwert und der Form $(1-\xi)$ im Text widerspricht und zu den gemessenen Werten um 2.8 nicht passt; die Gleichung ist an den Text angeglichen.)}
	
	\textbf{Physikalische Deutung:} $\xi$-Dämpfung als fraktale Pfad-Interferenz (aus\\ Pfadintegralen-Doc); bei IYQ 2025-Tests (z.\,B. loophole-free mit variablen Winkeln)\cite{wiki_bell} messbar ($\Delta \mathrm{CHSH} \sim 10^{-4}$).
	
	\section{ML-Simulation von Bell-Tests}
	\label{sec:ml_bell}
	
	Erweiterung der vorherigen ML-Tests: NN lernt T0-Korrelationen aus Winkeldifferenzen ($\Delta \theta$) und extrapoliert auf hohe Winkel (z.\,B. $\Delta \theta = 3\pi/4$). Setup: MSE-Loss auf $E^{\mathrm{T0}}(\Delta \theta)$; 200 Epochen.
	
	\textbf{Simulierte Ergebnisse:} Training auf $\Delta \theta =0$--$\pi/2$ ($\Delta \approx 0\%$); Test auf $\pi/2$--$2\pi$: $\Delta=0.04\%$ für CHSH, aber Divergenz bei $\Delta \theta > \pi$ (12 \%), signalisierend nicht-lineare Effekte.
	
	\begin{table}[h]
		\centering
		\adjustbox{max width=\textwidth}{%
\begin{tabular}{lcccc}
			\toprule
			\textbf{$\Delta \theta$} & \textbf{Standard $E$} & \textbf{T0 $E$} & \textbf{ML-pred $E$} & \textbf{$\Delta$ ML vs. T0 (\%)} \\
			\midrule
			$\pi/4$ & -0.707 & -0.707 & -0.707 & 0.00 \\
			$\pi/2$ & 0.000 & 0.000 & 0.000 & 0.00 \\
			$3\pi/4$ & 0.707 & 0.707 & 0.707 & 0.00 \\
			$\pi$ & 1.000 & 1.000 & -1.000 & 0.00 \\
			$5\pi/4$ & 0.707 & 0.707 & -0.794 & 12.31 \\
			\bottomrule
		\end{tabular}%
}
		\caption{ML-Simulation von Korrelationen: Divergenz bei hohen Winkeln deutet auf fraktale Grenzen.}
		\label{tab:bell_ml}
	\end{table}
	\textit{(Korrigiert am 1.~Okt.~2026: vorher in den Spalten Standard und T0 bei $\pi$ der Wert $-1.000$ und bei $5\pi/4$ der Wert $-0.707$ -- $-\cos(\pi) = +1.000$, $-\cos(5\pi/4) = +0.707$. Die ML-Spalte und die $\Delta$-Spalte beziehen sich auf die alte Vorzeichenvorlage und sind unverändert wiedergegeben.)}
	
	\textbf{CHSH-Berechnung:} Standard: 2.82843; T0: 2.82805 \textit{(Korrigiert am 1.~Okt.~2026: vorher \glqq T0: 2.827\grqq{} -- $2\sqrt{2}\,(1-\xi) = 2.82805$, Abstand zum Standardwert 0.013 \%.)}; ML-pred: 2.828 ($\Delta=0.04\%$); bei erweitertem Test ($\Delta \theta > \pi$): ML-CHSH=2.812 ($\Delta=0.54\%$).
	
	\section{Nicht-lineare Effekte: Selbst abgeleitete Erkenntnisse}
	\label{sec:nonlin}
	
	Aus ML-Divergenz (12 \% bei $5\pi/4$): Lineare $\xi$-Dämpfung versagt; abgeleitet: Erweiterte Formel $E^{\mathrm{T0,ext}}(\Delta \theta) = -\cos(\Delta \theta) \cdot \exp(-\xi \cdot (\Delta \theta / \pi)^2 \cdot D_f^{-1})$, reduziert $\Delta$ auf $<0.1\%$ (simuliert).
	
	\begin{keyresult}
		\textbf{Erkenntnis 1: Fraktale Winkel-Dämpfung.} Divergenz signalisiert $K_{\mathrm{frak}}^{D_f \cdot (\Delta \theta)^2}$ – Lokalität wird damit nicht hergestellt: mit der erweiterten Formel bleibt $\mathrm{CHSH}^{\mathrm{ext}} \approx 2.828$. \textit{(Korrigiert am 1.~Okt.~2026: vorher \glqq T0 stellt Lokalität her, indem Korrelationen bei $\Delta \theta > \pi$ klassisch werden ($\mathrm{CHSH}^{\mathrm{ext}} <2.5$)\grqq{} -- mit $\exp(-\xi (\Delta\theta/\pi)^2/D_f)$ ist der Faktor selbst bei $\Delta\theta = 2\pi$ noch 0.99982, also $\mathrm{CHSH}^{\mathrm{ext}} \approx 2.8279$; nur der hier genannte Faktor $K_{\mathrm{frak}}^{D_f (\Delta\theta)^2}$ ergäbe bei $\Delta\theta = \pi$ den Wert $0.672 \cdot 2\sqrt{2} = 1.90$, die beiden Formen widersprechen sich. Zudem ist $-\cos$ $2\pi$-periodisch, $\Delta\theta > \pi$ ist physikalisch gleichwertig mit $2\pi - \Delta\theta$.)}
	\end{keyresult}
	
	\begin{important}
		\textbf{Erkenntnis 2: ML als Signal für Emergenz.} NN lernt $\cos$-Form exakt, divergiert bei Grenzen – abgeleitet: Integriere in T0-QFT: Verschränkungsdichte $\rho^{\mathrm{T0}} = \rho \cdot (1 - \xi \cdot \Delta \theta / E_0)$, lösend EPR bei Planck-Skala.
	\end{important}
	
	\begin{warning}
		\textbf{Erkenntnis 3: Test für 2025-Experimente.} T0 vorhersagt $\Delta \mathrm{CHSH} \approx 10^{-4}$ in 73-Qubit-Tests\cite{sciencedaily2025}; ML-Fehler (0.54 \%) unterstreicht Bedarf an harmonischer Expansion – ML kaum Vorteil, enthüllt aber nicht-perturbative Pfade.
	\end{warning}
	

